\chapter{Réalisation}

\section{Fonctionnalités réalisées (Lot B)}
\begin{itemize}
  \item Développement des pages React/TypeScript pour les rôles enseignant et étudiant.
  \item Configuration du routage applicatif avec routes protégées et redirections par rôle.
  \item Intégration des modules : cours, devoirs, soumissions, notes, absences, demandes, réclamations, notifications.
  \item Mise en place de composants réutilisables et d'un écran générique de gestion des ressources.
  \item Validation des parcours fonctionnels de bout en bout côté interface.
\end{itemize}

\section{Résultats fonctionnels obtenus}
Le lot frontend a permis de livrer une interface exploitable avec des parcours métier complets. Les utilisateurs peuvent accéder à leurs fonctionnalités via un tableau de bord unifié et des pages dédiées.

Les points majeurs réalisés :
\begin{itemize}
  \item espace enseignant : gestion cours/devoirs/évaluations/emploi du temps ;
  \item espace étudiant : consultation cours, soumissions, notes, absences, demandes ;
  \item espace administrateur : supervision des ressources et modules de contrôle ;
  \item navigation cohérente et consommation fiable des APIs backend.
\end{itemize}

\section{Contraintes rencontrées et solutions apportées}
\begin{tabularx}{\textwidth}{lX}
\toprule
\textbf{Contrainte} & \textbf{Solution appliquée} \\
\midrule
Multiplicité des profils utilisateurs & Mise en place d'un routage protégé avec redirection par rôle. \\
Hétérogénéité des formulaires métier & Utilisation de composants réutilisables et d'une page générique ResourcePage. \\
Synchronisation avec les APIs backend & Alignement des champs frontend avec les contrats API et gestion explicite des erreurs. \\
Lisibilité de l'interface & Structuration par modules et amélioration de l'ergonomie de navigation. \\
\bottomrule
\end{tabularx}

\section{Compétences acquises}
Ce lot a permis de renforcer les compétences suivantes :
\begin{itemize}
  \item développement d'interfaces React/TypeScript ;
  \item conception de navigation applicative par rôles ;
  \item intégration d'API REST dans une SPA ;
  \item gestion des états UI (chargement, succès, erreur) ;
  \item validation fonctionnelle des parcours utilisateur.
\end{itemize}

\section{Livrables du Lot B}
\begin{itemize}
  \item pages frontend par rôle (admin, enseignant, étudiant) ;
  \item configuration complète des routes et menus de navigation ;
  \item composants réutilisables pour formulaires et ressources ;
  \item scénarios de test fonctionnels côté interface.
\end{itemize}

\section{Division des tâches du projet}
\begin{tabularx}{\textwidth}{lX}
\toprule
\textbf{Composant} & \textbf{Responsabilité Lot B (ce rapport)} \\
\midrule
Frontend & Architecture React, composants et routage applicatif. \\
Parcours utilisateurs & Implémentation des espaces enseignant/étudiant et formulaires métier. \\
Intégration API & Consommation des endpoints backend et traitement des états de chargement/erreur. \\
Qualité UX & Responsive design, cohérence visuelle, lisibilité et fluidité de navigation. \\
\bottomrule
\end{tabularx}
